% ===========================================================================
\subsection{\file{new-design-reference/js/main.js}}
% ===========================================================================

\textbf{JavaScript, 79 lines.} Everything on the page that reacts. It does three
things and nothing else: it makes the navigation bar grow a frosted background
once you scroll, it counts the four metric numbers up from zero when they come
into view, and it fades most blocks upward as they arrive.

\begin{provenance}{new-design-reference/js/main.js}
Hand-written, and its own header comment says as much on line 3: \code{Zero
external dependencies.} That claim is true and checkable: the file imports
nothing, requires nothing, and the repository has no \file{package.json} that
could supply a library.

It arrived complete in the first commit, \code{0d77d51}, and has never been
edited. The patterns in it (an immediately invoked function wrapper, an
\code{IntersectionObserver} factory, a cubic ease-out counter) are the standard
vocabulary of vanilla JavaScript rather than the fingerprint of any particular
source, so no upstream is claimed. \textbf{Confidence: certain for hand-written
and dependency-free; no provenance evidence found for any specific inspiration,
and none is invented here.}
\end{provenance}

\subsubsection*{Why it exists}

Without this file the page is complete, readable and static. Two of its three
jobs are pure decoration and would not be missed. The third is not: the metric
numbers are written into the HTML as the character \code{0}, with their real
values hidden in data attributes, so with this file gone the page permanently
advertises zero projects delivered and zero views.

\begin{keyidea}{The single most important line in this file is in another file}
\file{index.html} loads this script at line 595, immediately before
\code{</body>}. Everything below runs the instant it loads, with no wrapper
telling it to wait for the document. That only works because, by the time the
browser reaches line 595, every element already exists. Move the
\code{<script>} tag into the \code{<head>} and this file would run against an
empty page, find nothing, attach nothing, and report no error at all.
\end{keyidea}

\subsubsection*{Line by line}

\paragraph{Lines 1 to 6: the header, and the wrapper.}

\begin{lstlisting}[language=Java,firstnumber=1,caption={\file{js/main.js}, lines 1--6}]
/* ai.digitalise.agency · main.js
   Scroll-driven animations: nav gloss, counter triggers, section reveals.
   Zero external dependencies. */

(function () {
  'use strict';
\end{lstlisting}

Lines 1 to 3 are a comment, using the same \code{/* */} form as CSS. It is an
accurate summary of the file, which is worth noting because a comment that
describes what the code actually does is less common than it should be.

Line 5 opens a wrapper worth understanding, because it looks like noise and is
not.

\begin{jargon}{Immediately invoked function expression}
Writing \code{(function () \{ ... \})();} defines a function and runs it once,
straight away, without ever giving it a name. The point is not the running: it is
the boundary. Anything declared inside is invisible outside.

Without it, the names \code{nav}, \code{animateCounter} and \code{observe}
declared below would become global, meaning any other script on the page could
see them, and any other script that happened to use the same name would collide
with them. One stray global called \code{observe} is exactly the kind of bug that
takes an afternoon to find.
\end{jargon}

Line 6 is \code{'use strict';}, which switches the language into a stricter mode
for everything inside this wrapper.

\begin{jargon}{Strict mode}
A setting that makes JavaScript refuse things it would otherwise silently allow.
The most valuable one: without it, assigning to a name you never declared quietly
creates a global variable, so a typo such as \code{navv = 3} succeeds and does
nothing useful. In strict mode it throws an error and you find the typo
immediately.
\end{jargon}

\paragraph{Lines 8 to 16: the navigation bar.}

\begin{lstlisting}[language=Java,firstnumber=7,caption={\file{js/main.js}, lines 7--16}]

  /* ── Nav: add .scrolled glass on scroll ── */
  const nav = document.querySelector('.nav');
  if (nav) {
    const onScroll = () => {
      nav.classList.toggle('scrolled', window.scrollY > 40);
    };
    window.addEventListener('scroll', onScroll, { passive: true });
    onScroll();
  }
\end{lstlisting}

Line 9 finds the navigation bar. \code{document.querySelector} takes a CSS
selector, exactly the same syntax the stylesheet uses, and returns the first
matching element, or the special value \code{null} if there is none.

\begin{jargon}{const, and why not var}
\code{const} declares a name whose binding cannot later be pointed at something
else. It is the default choice in modern JavaScript, with \code{let} used when a
name genuinely must change. The older \code{var} has confusing scoping rules and
appears nowhere in this file, which is a small sign of care.
\end{jargon}

Line 10 is a guard: everything only happens \code{if (nav)}, that is, only if
something was actually found. This is the pattern that makes the whole file
tolerant of a changing page. Delete the navigation bar from the HTML and this
block is skipped rather than crashing.

\begin{jargon}{Truthy and falsy}
JavaScript treats certain values as false when tested in a condition:
\code{false}, \code{0}, an empty string, \code{null} and \code{undefined}.
Everything else is treated as true. So \code{if (nav)} reads as ``if an element
was found'', because \code{querySelector} returns \code{null} when it finds
nothing.
\end{jargon}

Lines 11 to 13 define what to do when the page scrolls, using an arrow function.

\begin{jargon}{Arrow function}
\code{() => \{ ... \}} is a compact way to write a function. Here the compactness
is the only benefit; the older \code{function () \{ ... \}} would behave
identically in this file.
\end{jargon}

Line 12 is the actual work, and it is one line doing the job of five.
\code{classList.toggle('scrolled', condition)} adds the class when the condition
is true and removes it when false. The condition is \code{window.scrollY > 40},
meaning the page has been scrolled more than 40 pixels from the top. The
stylesheet at lines 105 to 111 defines what \code{scrolled} looks like, and the
transition on line 103 makes the change fade rather than snap.

Line 14 registers the function to run on every scroll event. \code{\{ passive:
true \}} is a performance promise to the browser: this handler will never cancel
the scroll. That matters because without the promise, the browser must wait for
the handler to finish before it knows whether the scroll may proceed, which can
make scrolling feel sticky. It is the correct flag here and it is often forgotten.

Line 15 calls the function once immediately. This handles the case where the page
is loaded already scrolled, for example when a visitor refreshes halfway down or
arrives on a link ending in \code{\#book}. Without it the bar would be transparent
over content until the first scroll.

\paragraph{Lines 18 to 36: the counter.}

\begin{lstlisting}[language=Java,firstnumber=17,caption={\file{js/main.js}, lines 17--36}]

  /* ── Counter animation (cubic ease-out) ── */
  function animateCounter(el) {
    const raw    = el.dataset.target || '0';
    const suffix = el.dataset.suffix || '';
    const target = parseFloat(raw.replace(/[^0-9.]/g, ''));
    const isFloat = raw.includes('.');
    const dec     = isFloat ? (raw.split('.')[1] || '').length : 0;
    const dur     = 2200;
    const start   = performance.now();

    const step = (now) => {
      const p      = Math.min((now - start) / dur, 1);
      const eased  = 1 - Math.pow(1 - p, 3);
      const val    = eased * target;
      el.textContent = (isFloat ? val.toFixed(dec) : Math.round(val).toLocaleString()) + suffix;
      if (p < 1) requestAnimationFrame(step);
    };
    requestAnimationFrame(step);
  }
\end{lstlisting}

This is the most substantial piece of logic in the repository, and it is worth
taking slowly.

Line 20 reads the target. \code{el.dataset.target} is how JavaScript reads the
attribute \code{data-target} written in the HTML: strip the \code{data-} prefix
and what remains is a property of \code{dataset}. The \code{|| '0'} at the end
supplies a default: if the attribute is missing, \code{dataset.target} is
\code{undefined}, which is falsy, so the whole expression evaluates to
\code{'0'}.

Line 21 does the same for the suffix, defaulting to an empty string.

Line 22 turns the text into a number. \code{raw.replace(/[\^{}0-9.]/g, '')}
deletes every character that is not a digit or a dot, and \code{parseFloat} then
reads what is left as a number.

\begin{jargon}{Regular expression}
A pattern for matching text, written between slashes. \code{[0-9.]} means ``a
digit or a dot''; putting a caret first negates it, so it means ``anything that
is not a digit or a dot''. The \code{g} after the closing slash means global:
replace every occurrence, not just the first.
\end{jargon}

Line 23 asks whether the original text contained a dot, which is how the code
distinguishes \code{1.6} from \code{17}. Line 24 counts the digits after the dot,
so that \code{1.6} is displayed to one decimal place throughout the animation
rather than as \code{1.5999999}. \code{(raw.split('.')[1] || '')} takes everything
after the dot, falling back to an empty string if there is nothing there, and
\code{.length} counts it.

\begin{jargon}{Ternary operator}
\code{condition ? a : b} evaluates to \code{a} when the condition is true and
\code{b} otherwise. It is an \code{if} statement compressed into an expression,
and it appears twice in this file, on lines 24 and 32.
\end{jargon}

Line 25 fixes the duration at 2200 milliseconds, 2.2 seconds. Line 26 records the
start time using \code{performance.now()}, a clock designed for measuring
durations precisely.

Lines 28 to 34 define one frame of the animation, and this is the heart of it.

\code{const p = Math.min((now - start) / dur, 1)} computes progress: how far
through the animation we are, from 0 at the start to 1 at the end. The
\code{Math.min(..., 1)} caps it, so that if the browser is busy and a frame
arrives late, progress never exceeds 1 and the number never overshoots its target.
That cap is the difference between a counter that lands on 17 and one that
occasionally flashes 19 before settling.

\code{const eased = 1 - Math.pow(1 - p, 3)} is the easing. Linear progress feels
mechanical; this curve makes the count start fast and slow dramatically as it
approaches the target, which is what makes it feel like it is settling.

\begin{keyidea}{What the easing formula actually does}
Take progress \code{p}, subtract it from 1, cube the result, and subtract that
from 1. At \code{p = 0.5}, halfway through in time, the formula gives
$1 - 0.5^3 = 0.875$: the counter is already 87.5\% of the way to its target. The
whole second half of the animation covers the last eighth of the distance. That
asymmetry is the entire effect.
\end{keyidea}

Line 32 renders the number. If the target had a decimal point, \code{toFixed(dec)}
formats it to that many places. Otherwise \code{Math.round} makes it a whole
number and \code{toLocaleString()} inserts the thousands separators appropriate to
the visitor's region, so a large number reads as \code{1,200} for a British
visitor and \code{1.200} for a German one. The suffix is then glued on the end.

Line 33 asks for the next frame if the animation is not finished, and line 35
starts the first one.

\begin{jargon}{requestAnimationFrame}
A request to run a function just before the browser's next screen repaint, about
60 times a second. Calling it repeatedly from inside itself, as here, is the
standard way to run an animation: it stays in step with the display, and the
browser pauses it entirely when the tab is in the background, which a timer would
not.
\end{jargon}

\begin{honest}{The counter ignores the reduced-motion preference}
This function animates for 2.2 seconds regardless of whether the visitor has asked
their operating system to reduce motion. The stylesheet honours that preference
carefully for everything it controls (lines 762 to 768), so the intent is plainly
there; it simply was not carried into this file.

The fix is small and this document states it precisely so nobody has to design it:

\begin{lstlisting}[language=Java,caption={A four-line fix. This code is a
suggestion in this document and is NOT present in the repository.}]
const reduce = window.matchMedia('(prefers-reduced-motion: reduce)').matches;
if (reduce) {
  el.textContent = (isFloat ? target.toFixed(dec)
                            : target.toLocaleString()) + suffix;
  return;
}
\end{lstlisting}

Placed after line 26, that shows the final number immediately and skips the
animation entirely, which is the correct behaviour: the information is preserved
and only the movement is removed.
\end{honest}

\begin{honest}{The metric values are hidden from anything that does not run
JavaScript}
The real numbers live in \code{data-} attributes, and the visible text is
\code{0}. Search engines, link previews, and any visitor with JavaScript disabled
therefore see a page claiming zero of everything. Writing the true value in the
HTML and having the counter animate towards it from zero would remove that risk
entirely at no cost to the effect.
\end{honest}

\paragraph{Lines 38 to 53: the observer factory.}

\begin{lstlisting}[language=Java,firstnumber=37,caption={\file{js/main.js}, lines 37--53}]

  /* ── Intersection observer factory ── */
  function observe(selector, className, threshold, callback) {
    const els = document.querySelectorAll(selector);
    if (!els.length) return;

    const io = new IntersectionObserver((entries) => {
      entries.forEach((e) => {
        if (!e.isIntersecting) return;
        e.target.classList.add(className || 'visible');
        if (callback) callback(e.target);
        io.unobserve(e.target);
      });
    }, { threshold: threshold || 0.18 });

    els.forEach((el) => io.observe(el));
  }
\end{lstlisting}

This is the best-designed thing in the repository: a small general-purpose helper
that the four lines following it use to do four different jobs.

\begin{jargon}{IntersectionObserver}
A browser facility that watches elements and tells you when they enter or leave
the visible area. Before it existed, the only way to do this was to listen for
every scroll event and measure each element's position by hand, which ran hundreds
of times a second and made pages stutter. The observer does the work outside the
main thread and reports only when something actually changes.
\end{jargon}

Line 40 finds every matching element. \code{querySelectorAll} returns a list, which
may be empty, rather than \code{null}.

Line 41 guards against an empty list and returns early. \code{!els.length} reads as
``if the count is zero'', because \code{0} is falsy and \code{!} inverts it. This
is the same defensive habit as line 10.

Lines 43 to 50 create the observer. Its callback receives a list of entries, one
for each element whose visibility changed, and for each one it:

\begin{enumerate}
\item skips anything that is leaving rather than entering
  (\code{if (!e.isIntersecting) return});
\item adds the class, defaulting to \code{visible} when the caller did not name
  one;
\item runs the optional extra callback, which only the counter uses;
\item stops watching that element with \code{io.unobserve}.
\end{enumerate}

That fourth step is the one that matters most.

\begin{keyidea}{Why it stops watching}
Every reveal in this design is meant to happen once. Without \code{unobserve}, an
element that scrolled out of view and back would fire again, and for the metrics
that would mean the numbers restarting their count every time the visitor scrolled
past. Unobserving also releases the browser's bookkeeping for that element, so
what starts as thirty watched elements steadily becomes zero.
\end{keyidea}

Line 50 sets the threshold, defaulting to 0.18: the callback fires once 18\% of
the element is visible, rather than at the first pixel. Waiting for a fraction to
appear is what makes the reveal feel connected to the scroll instead of triggering
just off screen.

\begin{honest}{The default threshold is written in a way that cannot express zero}
\code{threshold || 0.18} substitutes the default whenever the argument is falsy,
and \code{0} is falsy. So a caller asking for a threshold of \code{0}, meaning
``fire the moment a single pixel appears'', would silently get \code{0.18}
instead. The same objection applies to \code{className || 'visible'}, though an
empty class name is not a sensible request.

No caller in this file passes \code{0}, so nothing is broken today. It is a latent
trap of a kind worth recognising, and the modern remedy is the \code{??} operator,
which substitutes only for \code{null} and \code{undefined} and leaves \code{0}
alone.
\end{honest}

Line 52 starts watching every element found.

\paragraph{Lines 55 to 67: the four things being watched.}

\begin{lstlisting}[language=Java,firstnumber=54,caption={\file{js/main.js}, lines 54--67}]

  /* reveal (fade + slide) */
  observe('.reveal', 'visible', 0.15);

  /* staggered children */
  observe('.stagger', 'visible', 0.15);

  /* pipeline animation */
  observe('.pipeline', 'visible', 0.2);

  /* counters */
  observe('.metric-num[data-target]', 'counted', 0.5, (el) => {
    animateCounter(el);
  });
\end{lstlisting}

Four calls, and the whole page's reveal behaviour. The first three add the class
\code{visible} to anything carrying \code{reveal}, \code{stagger} or
\code{pipeline}, at slightly different thresholds. The stylesheet does everything
after that.

The fourth is the different one. Its selector,
\code{.metric-num[data-target]}, matches elements that carry the class \emph{and}
have the attribute, which is a small piece of defensiveness: a metric element
without a target would be skipped rather than animated to \code{NaN}. Its
threshold is 0.5, so the numbers only start counting when the band is half visible,
and its callback runs the counter.

\begin{honest}{The class \code{counted} is added and styled nowhere}
The fourth call passes \code{'counted'} rather than \code{'visible'}, so each
metric number gains a class that no CSS rule and no other JavaScript ever reads. It
is harmless, and it serves as a marker a developer could inspect, but it is worth
knowing that it does nothing. Note also that the metric numbers are children of
\code{.metrics-grid}, which is a \code{.stagger}, so they are already being faded
in by the second call: the class here is genuinely spare.
\end{honest}

\paragraph{Lines 69 to 79: smooth anchor scrolling, and the close.}

\begin{lstlisting}[language=Java,firstnumber=68,caption={\file{js/main.js}, lines 68--79}]

  /* ── Smooth anchor scroll (book-a-call button) ── */
  document.querySelectorAll('a[href^="#"]').forEach((a) => {
    a.addEventListener('click', (e) => {
      const target = document.querySelector(a.getAttribute('href'));
      if (!target) return;
      e.preventDefault();
      target.scrollIntoView({ behavior: 'smooth', block: 'start' });
    });
  });

})();
\end{lstlisting}

Line 70 finds every link whose destination starts with a hash.
\code{a[href\^{}="\#"]} is an attribute selector, and \code{\^{}=} means ``starts
with''. On this page that matches three links: the navigation's ``Book a call'',
the hero's ``Book a call'', and the hero's ``See the work''.

Line 72 looks up the destination by using the link's own \code{href} as a
selector, which works because \code{\#book} is valid CSS meaning ``the element
with id book''.

Line 73 returns early if the destination does not exist, leaving the browser to
handle the click normally. Line 74 then cancels the browser's default jump, but
only after the destination has been confirmed, which is the correct order: cancel
first and a broken link would do nothing at all instead of at least trying.

Line 75 performs the scroll. \code{block: 'start'} aligns the destination with the
top of the window.

\begin{honest}{This whole block is redundant, and it carries a latent crash}
\file{css/site.css} line 28 already sets \code{scroll-behavior: smooth} on the
document, which makes every anchor link glide without any JavaScript at all. These
eleven lines reproduce that.

They also introduce a failure the CSS version does not have. A link written as
\code{href="\#"}, which is a common way to write a placeholder link, would reach
line 72 as \code{document.querySelector('\#')}. That is not a valid selector, so
the browser throws a \code{SyntaxError}, and because the error happens inside a
click handler it would break that click and leave an error in the console.

No such link exists in the page today, so this is latent rather than live. The
honest summary is that the safest change to this file would be to delete lines 69
to 77 entirely: it would remove a redundancy and a hazard in one edit, and nothing
visible would change.
\end{honest}

Line 79 closes and immediately runs the wrapper opened on line 5.

\subsubsection*{What it connects to}

\textbf{Named by:} \file{new-design-reference/index.html}, line 595, and nothing
else.

\textbf{Names:} no file at all. It imports nothing and fetches nothing.

\textbf{Its entire interface with the rest of the project is seven strings}, and
they are worth listing because they are the contract nothing checks:

\begin{table}[H]\centering\small
\begin{tabular}{@{}lll@{}}
\toprule
String & Used at & Where it must also exist \\
\midrule
\code{.nav}        & line 9  & \file{index.html} line 23, \file{site.css} line 95 \\
\code{scrolled}    & line 12 & \file{site.css} line 105 \\
\code{.reveal}     & line 56 & \file{index.html} (five places), \file{site.css} line 78 \\
\code{.stagger}    & line 59 & \file{index.html} (four places), \file{site.css} line 86 \\
\code{.pipeline}   & line 62 & \file{index.html} line 223, \file{site.css} line 381 \\
\code{visible}     & lines 56--62 & \file{site.css} lines 83, 92, 400 \\
\code{.metric-num} & line 65 & \file{index.html} (four places), \file{site.css} line 343 \\
\bottomrule
\end{tabular}
\caption{Every point of contact between this file and the other two. Break any
one of these strings in one place and the page keeps working, minus one
behaviour, with no error reported anywhere.}
\end{table}

\subsubsection*{How it breaks}

\begin{itemize}
\item \textbf{Load it in the \code{<head>} instead of at the end of the body} and
  all three behaviours vanish silently, since every lookup runs immediately and
  finds nothing.
\item \textbf{Rename a class in the HTML} and the corresponding behaviour stops,
  again silently, because of the early-return guards.
\item \textbf{Delete this file} and the page is static, with four zeros in the
  metrics band. That last part is the only real damage.
\item \textbf{Add a link with \code{href="\#"}} and clicking it throws a
  \code{SyntaxError} from line 72.
\item \textbf{A visitor who has asked for reduced motion} still gets the counting
  animation, because this file never checks.
\item \textbf{An ancient browser without \code{IntersectionObserver}} would throw
  on line 43 and lose all four reveal behaviours, which would leave every
  \code{.reveal} element permanently invisible were it not for the CSS
  reduced-motion block. Every browser in current use supports it, so this is a
  historical note, not a live concern.
\end{itemize}

\subsubsection*{Honest findings for this file}

The counter ignores \code{prefers-reduced-motion}; the real metric values are
invisible without JavaScript; the smooth-scroll block duplicates a CSS feature and
carries a latent \code{SyntaxError}; the \code{counted} class is inert; and the
\code{||} defaults on line 47 and 50 cannot express a falsy argument. Against
those: the file is genuinely well written, with consistent guards, a sensible
shared helper, correct use of \code{unobserve} and \code{passive}, and no
dependencies at all.
